\documentclass[main.tex]{subfiles}

\begin{document}
\section{Analisi e compilazione del codice vulnerabile}\label{sec:cap 2}

Per sviluppare il nostro attacco di buffer overflow abbiamo innanzitutto creato un programma vulnerabile in c, chiamato vuln.c.\\

\begin{lstlisting}[style=codiceC]
#include <stdio.h>
#include <string.h>

void vulnerable() {
    char buffer[100];
    printf("Inserisci input: ");
    gets(buffer);  
}

int main() {
    vulnerable();
}
\end{lstlisting}

\noindent La funzione \texttt{gets()} è intrinsecamente insicura poiché legge l'input da standard input (\texttt{stdin}) senza controllare la lunghezza della stringa inserita. Questo può causare una \textbf{scrittura oltre i limiti dell'array} \texttt{buffer[100]} se l'utente inserisce più di 99 caratteri (più il terminatore \texttt{\textbackslash 0}).

\subsubsection*{Compilazione del codice e disattivazione dei sistemi di sicurezza}

Per poter analizzare e sfruttare la vulnerabilità di tipo buffer overflow, è necessario compilare il programma disattivando alcune protezioni.\\
Innanzitutto, compiliamo il programma vulnerabile vuln.c, generando l'eseguibile, nel seguente modo:
\begin{verbatim}
gcc -m32 -fno-stack-protector -z execstack -no-pie -g vuln.c -o vuln
\end{verbatim}

\noindent Significato delle opzioni

\begin{itemize}
    \item \verb|-m32|: compila in modalità 32 bit.
    \item \verb|-fno-stack-protector|: disattiva lo \textbf{Stack Canary}, una protezione che rileva sovrascritture dello stack e può prevenire exploit.
    \item \verb|-z execstack|: rende lo \textbf{stack eseguibile}, permettendo di eseguire codice inserito nello stack (la shellcode nel nostro caso).
    \item \verb|-no-pie|: disattiva la compilazione come \textbf{Position Independent Executable (PIE)}. Questo rende gli indirizzi di memoria fissi, facilitando l'attacco.
    \item \verb|-g|:  permette di effettuare il debug del programma

\end{itemize}

\noindent Successivamente dobbiamo disattivare l'\textbf{ASLR} (Address Space Layout Randomization),  una protezione del sistema operativo che randomizza gli indirizzi di memoria ad ogni esecuzione, nel seguente modo:
\begin{verbatim}
echo 0 | sudo tee /proc/sys/kernel/randomize_va_space
\end{verbatim}
per controllare se tutte le sicurezze sono state disattivate utilizziamo il comando da terminale \verb|checksec --file=./vuln| aspettandoci un output di questo tipo:

\begin{itemize}
    \item Canary found:          No
    \item NX enabled:            No
    \item PIE enabled:           No
    \item Stack executable:      Yes
\end{itemize}












\end{document}